% secciones/modelo.tex -- Sec 2.2 Modelo

El proceso se describe con un modelo dinámico de parámetros concentrados en equivalentes de carbono, que toma como núcleo el modelo de asimilación, crecimiento y producción de CO$_2$ de \textcite{eriksenDynamicModellingFeed2022} para \textit{H. illucens} y añade una extensión: una fuente microbiana asociada con la degradación del alimento no consumido. El estado integrado es $\mathbf{y} = (B, L_{\mathrm{tot}}, D_C, X)$, donde $B$ y $L_{\mathrm{tot}}$ son la biomasa estructural y los lípidos totales de la larva (mg~C larva$^{-1}$), $D_C$ el carbono del alimento disponible en la cámara (g~C) y $X$ la biomasa microbiana del lecho (g~C); el carbono respirado acumulado se lleva como variable contable. El modelo supone temperatura constante (medida), oxígeno no limitante dentro de los cierres y número de larvas $N$ constante. El origen temporal $t = 0$ es el montaje del experimento, con larvas de 4.5~d de edad; el núcleo de Eriksen se gobierna por la edad fisiológica $\tau = t + 4.5$~d, de modo que el umbral $t_{p,\min} = 13$~d desde la eclosión equivale a 8.5~d desde el montaje.

La cámara de medición mezcla dos fuentes de CO$_2$ que el sensor no puede separar: la respiración de las larvas y la degradación microbiana del alimento. Atribuir el total medido a cada fuente exige un término mecanístico para ambas. La larva se modela como un procesador de carbono: ingiere alimento y asimila una fracción proporcional a su tamaño (tasa $r_A$); ese carbono se reparte entre crecimiento de la estructura corporal ($B$), reservas lipídicas ($L_{\mathrm{tot}}$) y costos metabólicos pagados como CO$_2$ (mantenimiento $r_{C,m}$, síntesis de estructura y depósito de lípidos). Así, la emisión larvaria es la suma de esos costos: todo carbono asimilado que no permanece en el cuerpo sale como gas. En paralelo, el alimento no consumido se degrada en el lecho por la microbiota ambiental, con cinética propia. El resto de esta sección presenta el núcleo que describe la larva \parencite{eriksenDynamicModellingFeed2022}, después la fuente microbiana del lecho y al final el observable de cámara con el que se compara el modelo. El modelo de referencia no representa la fuente microbiana del lecho; se introduce en este artículo.

\paragraph{Estadios de desarrollo.}
Las larvas de \textit{H. illucens} crecen por mudas: entre muda y muda (un instar) su biomasa corporal se multiplica por un factor aproximadamente constante. El modelo no describe cada muda: comprime el desarrollo en tres estadios operativos según la biomasa estructural $B$ (el tejido corporal, sin las reservas lipídicas), con dos referencias de tamaño: la biomasa máxima $B_{\max,0}$ (el mayor tamaño que puede alcanzar una larva sin restricciones) y el umbral $B_{\max,0}/\delta_I$, donde $\delta_I = 5$ es el incremento relativo de peso por muda. Mientras $B < B_{\max,0}/\delta_I$ la larva está en el estadio inicial ($I = 5$), de asimilación plena; cuando $B_{\max,0}/\delta_I \leq B < B_{\max}(\tau)$ ha entrado en el último estadio alimentario ($I = 6$), con la asimilación frenada; y al alcanzar $B \geq B_{\max}(\tau)$ ---el techo de biomasa a la edad $\tau$--- pasa a prepupa ($I = 7$), que ya no se alimenta. La numeración no es arbitraria: el modelo de referencia inicia su descripción en el instar 5, el estadio de sus larvas iniciales (~5~d de edad), y numera los estadios consecutivamente hasta la prepupa; se conserva esa numeración para que ambos modelos sean directamente comparables.

El máximo de biomasa que una larva puede alcanzar no es fijo. En la especie, el peso final y el tiempo de desarrollo se relacionan negativamente: las larvas que tardan más en crecer pupan a menor tamaño \parencite{eriksenDynamicModellingFeed2022}. El modelo lo representa con un techo que decrece con la edad a partir de un tiempo mínimo de desarrollo:

\begin{equation}
  B_{\max}^{\text{crudo}}(\tau) = B_{\max,0} - \rho \max(\tau - t_{p,\min}, 0),
  \label{eq:techo-crudo}
\end{equation}
y aplica después la protección
\begin{equation}
  B_{\max}(\tau) = \max!\left(B_{\max}^{\text{crudo}}(\tau), \varepsilon_B\right),
  \label{eq:techo}
\end{equation}
con $t_{p,\min} = 13$~d (ni la larva más rápida pupa antes) y $\varepsilon_B = 10^{-12}$~mg~C. El valor de referencia es $\rho = 1$~mg~C~d$^{-1}$, la razón a la que se estrecha la meta para las larvas lentas; en C este parámetro se recalibra. El piso no representa un tamaño fisiológico: evita techos nulos o negativos y protege los denominadores de las funciones de asimilación y crecimiento. El valor protegido es también el que entra en la clasificación de instares. $B_{\max,0} = 39.69$~mg~C se hereda del modelo de referencia: es su valor por defecto de 90~mg de peso seco (dentro del rango 65--90~mg reportado por \citeauthor{eriksenDynamicModellingFeed2022}), escalado a carbono orgánico multiplicando por $\delta_{C,B} = 0.49$ (fracción de carbono de la biomasa estructural) y por $\delta_O = 0.9$ (fracción orgánica de esa biomasa; el resto es ceniza). En este estudio no se reajusta (sec.~\ref{sec:calibracion}). Estos coeficientes de composición, junto con $\delta_{C,L}$ y $\delta_{L,B}$, se emplean en todo el modelo para convertir entre peso fresco, peso seco y carbono; su lista completa está en el Material Suplementario (sec.~\ref{sup:params}).

\paragraph{Adquisición y destino del carbono.}
Fijado el programa de desarrollo, el modelo describe cómo entra el carbono a la larva y a dónde va. La asimilación es proporcional al tamaño del cuerpo. La intensidad específica de asimilación $a$ depende del estadio:
\begin{equation}
  r_A = a\,B,
  \label{eq:asim}
\end{equation}
con
\begin{equation}
  a = \begin{cases}
    a_{\max}, & I = 5, \\[4pt]
    a_{\max}\left(1 - (B/B_{\max}(\tau))^{\alpha}\right), & I = 6, \\[4pt]
    0, & I = 7,
  \end{cases}
  \label{eq:a-estadio}
\end{equation}
donde $a_{\max} = 1$~d$^{-1}$ es la tasa específica máxima de asimilación ---el techo fisiológico de cuánto carbono incorpora la larva por unidad de biomasa estructural y por día--- y $\alpha = 1$. En el estadio 5 la asimilación es plena; en el 6 el factor de Verhulst la frena a medida que $B$ se acerca al techo; en prepupa ($I = 7$) cesa por completo.

La calidad de la dieta modula la asimilación mediante un factor adimensional $f_d$:
\begin{equation}
  a_{\mathrm{ef}} = a\,f_d,
  \label{eq:aef}
\end{equation}
con $f_d = 1$ en D1 (la dieta con que se calibró el modelo de referencia) y el factor de D4 estimado de los datos de este experimento (sec.~\ref{sec:calibracion}). La forma multiplicativa se elige por economía: es la modificación más leve posible, que no toca las ecuaciones, los parámetros ni el programa de desarrollo del núcleo, y solo reduce la intensidad con que la larva asimila.

El carbono asimilado se reparte con prioridad fija. Primero, el mantenimiento:
\begin{equation}
  r_{C,m} = m\,B,
  \label{eq:mant}
\end{equation}
proporcional a la estructura, con prioridad absoluta. Segundo, el crecimiento:
\begin{equation}
  r_B = \mu_B\,B,
  \label{eq:crec}
\end{equation}
con tasa específica
\begin{equation}
  \mu_B = \frac{(a_{\mathrm{ef}} - m)}{1 + Y_B}\left(1 - \left(\frac{B}{B_{\max}(\tau)}\right)^{\beta}\right),
  \label{eq:muB}
\end{equation}
y $\beta = 1$. Los tres factores se leen por separado: $(a_{\mathrm{ef}} - m)$ es lo que sobra de la asimilación después de pagar el mantenimiento; $1/(1 + Y_B)$ convierte ese remanente en crecimiento neto, porque fabricar estructura cuesta ---$Y_B$ unidades de CO$_2$ por cada unidad depositada: con $Y_B = 0.44$, de cada 1.44 unidades asignadas al crecimiento, 1 queda en el cuerpo y 0.44 se respiran---; y el factor logístico frena el crecimiento a medida que la larva se acerca a su techo de tamaño. En prepupa el crecimiento cesa, $r_B = 0$.

Tercero, las reservas. El excedente después de mantenimiento y crecimiento es
\begin{equation}
  s_L = r_A - (r_B + r_{C,m} + Y_B\,r_B).
  \label{eq:excedente}
\end{equation}
Si $s_L \geq 0$, el remanente se deposita como lípido:
\begin{equation}
  r_L = \frac{s_L}{1 + Y_L}, \qquad r_{C,L} = Y_L\,r_L,
  \label{eq:lipidos-pos}
\end{equation}
donde $Y_L$ es el costo de depósito ---la versión de $Y_B$ para el almacén--- y $r_{C,L}$ el CO$_2$ asociado. Si $s_L < 0$ ---la asimilación no alcanza a pagar siquiera el mantenimiento, como en ayuno--- no hay nada que depositar: la larva moviliza sus reservas a razón
\begin{equation}
  r_L = s_L,
  \label{eq:lipidos-neg}
\end{equation}
un flujo negativo que drena $L_{\mathrm{tot}}$, sin costo de depósito porque desmantelar el almacén no cuesta lo mismo que llenarlo.

Los balances de estado recogen el resultado del reparto:
\begin{align}
  \frac{dB}{d\tau} &= r_B, \label{eq:balanceB} \\
  \frac{dL_{\mathrm{tot}}}{d\tau} &= r_L + \delta_{L,B}\,r_B,
  \label{eq:balanceL}
\end{align}
con $\delta_{L,B} = 0.1$: el tejido corporal no es proteína pura e incorpora alrededor de un 10% de lípido al crecer.

\paragraph{El reparto no deja residuos.}
Todo carbono que entra por $r_A$ termina en el cuerpo ($B$), en el almacén ($L_{\mathrm{tot}}$) o en el aire como CO$_2$ (los tres costos): el reparto no deja residuos. Esta cuenta cerrada es la base de la separación de fuentes que persigue este estudio: si el destino de cada unidad asimilada está resuelto por el modelo, la fracción que respira la larva queda determinada y no debe adjudicarse al total de la cámara.

\paragraph{La emisión larvaria se calcula, no se mide.}
La emisión de CO$_2$ de la larva no es un dato experimental: se obtiene sumando sus costos respiratorios:
\begin{equation}
  r_{CO_2,\mathrm{lar}} = r_{C,m} + r_{C,B} + r_{C,L} + r_{\mathrm{cat}},
  \label{eq:emision}
\end{equation}
con $r_{C,B} = Y_B\,r_B$ el costo de síntesis estructural, $r_{C,L} = Y_L\,\max(r_L, 0)$ el costo de depósito lipídico ---solo se cobra cuando hay síntesis: movilizar reservas en ayuno no cuesta--- y $r_{\mathrm{cat}}$ el término catabólico del cierre del ciclo, nulo fuera de prepupa. En masa de CO$_2$, la tasa se multiplica por $44/12$. Esta emisión calculada es el aporte larvario que, sumado al microbiano, se compara con el sensor de la cámara en los resultados.

Las condiciones iniciales traducen el tamaño de las larvas al montaje a las unidades del modelo. La larva se siembra con un peso fresco medido en el ensayo de crecimiento; el peso seco $X_{DW,0}$ (mg) se obtiene del peso fresco mediante la fracción de materia seca. De ese peso seco solo una parte es carbono en la unidad del modelo: la fracción orgánica $\delta_O = 0.9$ (el 10% restante del peso seco es ceniza) contiene a su vez una fracción de carbono $\delta_{C,B} = 0.49$. Estos coeficientes de composición ---como $\delta_{L,B}$ y $\delta_{C,L}$--- se toman de la Tabla~1 del modelo de referencia \parencite{eriksenDynamicModellingFeed2022}: describen la composición química típica de la larva y no se miden en este estudio. El producto $B_0 = X_{DW,0}\,\delta_{C,B}\delta_O$ es el carbono que la larva trae en su cuerpo al montaje. Los lípidos iniciales se estiman con la misma lógica: una fracción $\delta_{L,B} = 0.1$ de la biomasa estructural es lípido (los lípidos del propio tejido), y ese lípido es carbono en una fracción $\delta_{C,L} = 0.79$ (valor del lípido de referencia del modelo, la trilaurina). El convenio del modelo de referencia asume que las reservas iniciales son exactamente esos lípidos estructurales: $L_{\mathrm{tot},0} = B_0\,\delta_{L,B}\,\delta_{C,L}/(\delta_{C,B}\delta_O)$. Qué parámetros se heredan, cuáles se fijan y cuáles se calibran se define en la sec.~\ref{sec:calibracion}; la procedencia valor por valor está en el Material Suplementario (sec.~\ref{sup:params}).

\paragraph{El cierre del ciclo: prepupa.}
Entrar en prepupa no congela a la larva: deja de asimilar ($r_A = 0$), pero sigue respirando y empieza a consumir su propia estructura. El mantenimiento baja a un mínimo de supervivencia:
\begin{equation}
  r_{C,m} = f_{\mathrm{mant}}\,m\,B,
  \label{eq:prep-mant}
\end{equation}
con $f_{\mathrm{mant}} \in (0,1]$ un factor de reducción que se calibra; en paralelo, la estructura se cataboliza a tasa proporcional:
\begin{equation}
  \frac{dB}{d\tau} = -k_{\mathrm{cat}}\,B,
  \label{eq:prep-cat}
\end{equation}
y el carbono catabolizado se respira entero:
\begin{equation}
  r_{\mathrm{cat}} = k_{\mathrm{cat}}\,B
  \quad \text{(mg~C larva$^{-1}$ d$^{-1}$)},
  \label{eq:rcat}
\end{equation}
sin pasar por el excedente de lípidos. Como $r_A = 0$, el excedente es $s_L = -r_{C,m} < 0$: las reservas pagan el mantenimiento reducido movilizándose a tasa $r_L = s_L$, sin costo de depósito, y no hay costos de síntesis ($r_{C,B} = r_{C,L} = 0$), de modo que la emisión de la prepupa es $r_{CO_2,\mathrm{lar}} = r_{C,m} + r_{\mathrm{cat}}$. La cuenta cierra exacto: como $L_{\mathrm{tot}}$ incluye los lípidos estructurales, el carbono corporal total se contabiliza sin duplicarlos como $C_{\mathrm{cuerpo}} = (1-\delta_{L,B})B + L_{\mathrm{tot}}$; con el balance anterior, en prepupa $dC_{\mathrm{cuerpo}}/d\tau = -(k_{\mathrm{cat}}B + f_{\mathrm{mant}}mB)$, que es exactamente lo respirado. La conservación del carbono vale entonces en los tres estadios. El módulo es además una extensión estricta del núcleo: con $f_{\mathrm{mant}} = 1$ y $k_{\mathrm{cat}} = 0$ reproduce exactamente la prepupa original, inmóvil y pagando $mB$ con lípidos. La clasificación se reevalúa sin histéresis; por debajo del techo, la regla asigna el instar 6 solo si $B \geq B_{\max,0}/\delta_I$. Los dos parámetros se estiman en la etapa 3 (sec.~\ref{sec:calibracion}); la estimación de $f_{\mathrm{mant}}$ alcanzó la frontera inferior de su intervalo de búsqueda.

\paragraph{El lecho microbiano: una segunda fuente de CO$_2$.}
Las larvas no consumen todo el sustrato: una fracción permanece en el lecho, donde la microbiota ambiental la descompone en paralelo al consumo larvario. Los controles sin larvas del experimento lo confirman: producen CO$_2$ de forma sostenida, señal que no puede atribuirse a la respiración larvaria. Esta fuente ---ausente en el modelo de referencia--- se describe con dos variables ya presentadas, el carbono disponible $D_C$ y la biomasa microbiana $X$ (g C), y una cinética de degradación proporcional a la biomasa microbiana y limitada por el sustrato:
\begin{equation}
  r_{\mathrm{deg}} = k\,X\,\frac{D_C}{K_{SM} + D_C}
  \quad \text{(g C d$^{-1}$)},
  \label{eq:deg}
\end{equation}
con $K_{SM}$ la constante de semisaturación por sustrato. El carbono degradado se reparte en tres vías: una fracción $Y_X$ crea biomasa microbiana nueva, una fracción $Y_{C,\mathrm{mic}}$ se emite como CO$_2$ ($r_{C,\mathrm{mic}} = Y_{C,\mathrm{mic}}\,r_{\mathrm{deg}}$) ---este es el aporte microbiano de la cámara--- y el resto permanece como residuo estabilizado.

La constante efectiva $k$ integra los dos factores ambientales que controlan la actividad microbiana documentados para este proceso \parencite{bekkerImpactSubstrateMoisture2021a}: la temperatura, con corrección tipo van 't Hoff de factor $Q_{10}$ sobre 25~°C, y la humedad del lecho, con respuesta tipo Monod de semisaturación $K_{\theta}$:
\begin{equation}
  k = k_{\mathrm{ref}}\,Q_{10}^{(T-25)/10}\,\frac{\theta}{K_{\theta} + \theta},
  \label{eq:k-microb}
\end{equation}
Como el experimento no registra la evolución de la humedad, $\theta$ se trata como parámetro fijo del ensayo (valor en el Material Suplementario). El balance del alimento cierra el módulo:
\begin{equation}
  \frac{dD_C}{dt} = -\frac{N\,r_A}{1000} - r_{\mathrm{deg}},
  \label{eq:balanceDC}
\end{equation}
con las cargas de alimento (500~g al montaje, 250~g al día~5) incorporadas como pulsos instantáneos convertidos a carbono. Hay una limitación de identificabilidad: los controles sin larvas informan el producto $k\,X$ pero no permiten separar $k$, $X$ y $Y_X$; se fijan entonces $X_0 = 1$~g C y $Y_X = 0.4$ (rendimiento de crecimiento aerobio típico \parencite{roels1980,shulerbioprocess2017}) y se calibran $k_{\mathrm{ref}}$ y la fracción de carbono degradable del alimento $\delta_{CS}$ contra las curvas de los controles (sec.~\ref{sec:calibracion}). Se declara un límite del modelo: la asimilación larvaria $r_A = a_{\mathrm{ef}}B$ permanece positiva aunque $D_C$ caiga a cero. En este experimento el alimento nunca se agota ---las cargas totales superan ampliamente el consumo---, de modo que el modelo opera dentro de su rango de validez; fuera de él, cuando $D_C \to 0$, la formulación se vuelve mecánicamente inconsistente porque la larva seguiría asimilando sin sustrato.

La biomasa microbiana tampoco se acumula indefinidamente: una fracción se autodestruye y respira su propio carbono, en la tradición de los modelos de fase endógena del crecimiento microbiano \parencite{roels1980}. El modelo lo representa con un decaimiento de primer orden, $k_d\,X$, que se resta del crecimiento y cuyo carbono se emite como CO$_2$ microbiano:
\begin{equation}
  \frac{dX}{dt} = Y_X\,r_{\mathrm{deg}} - k_d\,X - G,
  \label{eq:balanceX}
\end{equation}
donde $k_d = 0.05$~d$^{-1}$ es la tasa de decaimiento, fija de la literatura de heterótrofos aerobios (tiempo de recambio $\approx 20$~d; sec.~\ref{sup:params}), y $G$ es el pastoreo larvario. A medida que el sustrato se agota, $r_{\mathrm{deg}}$ cae y el decaimiento hace que $X$ alcance un máximo y decrezca.

En los tratamientos existe un segundo control sobre la biomasa microbiana, que solo opera en presencia de larvas: el pastoreo. Las larvas de \textit{H. illucens} ingieren microbios junto con el sustrato; el modelo lo representa con un flujo proporcional a la asimilación larvaria y saturado en la biomasa disponible:
\begin{equation}
  G = \gamma\,\frac{N\,r_A}{1000}\,\frac{X}{K_G + X},
  \label{eq:pastoreo}
\end{equation}
con $\gamma$ el coeficiente de pastoreo (gramos de carbono microbiano pastoreados por gramo de carbono asimilado por la larva) y $K_G = 0.1$~g~C la semisaturación del encuentro. El carbono pastoreado se incorpora al residuo: a esta resolución del modelo, la larva empaca la biomasa ingerida en el frass sin asimilarla. Dos propiedades hacen al término bien comportado: sin larvas no hay pastoreo ($G = 0$ en los controles) y en ayuno o prepupa se apaga solo ($r_A = 0$), de modo que el florecimiento microbiano queda acotado solo donde las larvas están activas. $\gamma$ se calibra (sec.~\ref{sec:calibracion}); su estimación interior al intervalo de búsqueda confirma que el control top-down es identificable con estos datos.

\paragraph{Del modelo al sensor: el observable de cámara.}
El modelo produce tasas de carbono; el sensor registra una pendiente de concentración en ppm/min, y comparar ambos exige traducir una a la otra. El carbono respirado por la cámara completa es
\begin{equation}
  R_{C,\mathrm{total}} = \frac{N\,r_{CO_2,\mathrm{lar}}}{1000} + r_{C,\mathrm{mic}}
  \quad \text{(g C d$^{-1}$)},
  \label{eq:observable}
\end{equation}
que se convierte en moles ($\div 12$) y, en la cámara estática bien mezclada, la pendiente de concentración sigue la ley del gas ideal:
\begin{equation}
  \frac{dC_{\mathrm{ppm}}}{dt} = 10^{6}\,\frac{R_u\,T}{P\,V_{\mathrm{aire}}}\,\frac{R_{C,\mathrm{total}}}{12 \times 1440}
  \quad \text{(ppm min$^{-1}$)},
  \label{eq:gas-ideal}
\end{equation}
con $V_{\mathrm{aire}}$ el volumen de aire efectivo de la cámara, $T$ la temperatura absoluta y $P$ la presión atmosférica. Entre cierres la cámara se ventila, de modo que la concentración no se acumula de un día a otro: el observable del modelo es la pendiente instantánea evaluada en los días de medición, no una trayectoria de concentración.

\paragraph{El techo del sensor: cotas, no puntos.}
En la mitad de los puntos de tratamiento los cierres alcanzaron la saturación del sensor; allí la lectura deja de seguir a la tasa real y el ajuste lineal subestima la emisión. Esos puntos no se descartan ni se tratan como valores exactos: se usan como cotas inferiores ---el modelo solo incurre en residual si su predicción queda por debajo de lo observado---, criterio que se formaliza con la calibración (sec.~\ref{sec:calibracion}). La fracción de puntos censurados y su efecto sobre los resultados se reportan en la sección de resultados.

\paragraph{Fidelidad de la implementación.}
La implementación de este estudio se verificó contra el archivo de modelación que \citeauthor{eriksenDynamicModellingFeed2022} publica como material suplementario, con comparación fórmula a fórmula. La verificación dejó dos divergencias documentadas entre las ecuaciones impresas del modelo de referencia y su propio archivo: el umbral de prepupa impreso usa el techo constante $B_{\max,0}$ mientras el archivo usa el techo decreciente $B_{\max}(\tau)$, y el valor impreso $\delta_{C,L} = 0.74$ aparece como $0.79$ en el archivo. Esta implementación reproduce el archivo en ambos casos. No existe la divergencia adicional de un máximo en cero en la ecuación de crecimiento: el texto y el archivo emplean la misma regla. De la misma manera, el reparto instantáneo descrito en esta sección reproduce el comportamiento del compartimento transitorio de asimilados del modelo impreso, que su propio archivo tampoco integra. La tabla de correspondencia completa (ecuación impresa $\leftrightarrow$ archivo $\leftrightarrow$ código de este estudio), junto con la comparación de métodos de integración, se reporta en el Material Suplementario (sec.~\ref{sup:correspondencia}).
